% Consolidated textbook-ch08 from f43cd09215a13d29f88baa670fa48441ec2a44e8; source fragments preserved.

% === BEGIN textbook-ch08-choices.tex blob 342dc34e310b0fa53df74b7e36c3081e8b551389 ===
% 《大学物理简明教程》第5版，第8章稳恒磁场；选择题8.1.1—8.1.13。
% 用户扫描PDF/印刷216—218页，高分辨率逐题核验；题意研究性改述，答案独立推导。
\subsection{选择题：毕奥--萨伐尔定律、安培环路定理与磁介质}

\sourceproblem{8.1.1}{216}{216}
\problemtext 真空中半径$R$的半圆形细导线流有稳恒电流$I$，求圆心$O$的磁感应强度大小。
\workings
圆弧张角$\alpha=\pi$，由毕奥--萨伐尔定律积分
\[
 B_O=\frac{\mu_0 I}{4\pi R}\int_0^\pi d\varphi
 =\boxed{\frac{\mu_0 I}{4R}}.
\]
原书选\textbf{D}。方向由右手法则判定，题只问大小。
\crosscheck 整圆的中心磁场为$\mu_0I/(2R)$，半圆恰为其一半。

\sourceproblem{8.1.2}{216}{216}
\problemtext 关于安培环路定理：若环流$\oint_L\bm H\cdot d\bm l=0$，
是否意味着曲线上每点$\bm H=0$、一定无电流穿过、或所围穿导电流的代数和为零？
\workings
稳恒电流下
\[
\oint_L\bm H\cdot d\bm l=I_{\mathrm f,\mathrm{enc}}^{\rm alg}.
\]
环流为零只表示按曲线定向计算的\emph{自由传导电流代数和}为零，
不排除大小相等、方向相反的穿越电流，也不排除外部电流在曲线上产生非零磁场。
故原书选\textbf{C}。
\crosscheck 在真空中$\bm B=\mu_0\bm H$，
环流积分为零不意味着被积向量逐点为零。

\sourceproblem{8.1.3}{216}{216}
\problemtext 半径$R$的无限长实心圆柱体内均匀分布轴向稳恒电流。
判断磁感应强度$B(r)$随距轴线$r$的变化曲线（原书图8.1.3）。
\workings
横截面均匀电流密度$J=I/(\pi R^2)$。
安培定理给
\[
\boxed{B(r)=
\begin{cases}
\mu_0Ir/(2\pi R^2),&0\le r<R,\\
\mu_0I/(2\pi r),&r\ge R.
\end{cases}}
\]
所以中心$B(0)=0$，内部随$r$线性增长，表面达到峰值，
外部按$1/r$递减且在$r=R$处连续。按原图选\textbf{A}。
\crosscheck 这是\emph{均匀体电流}，不是所有电流集中在薄圆筒表面的模型。

\sourceproblem{8.1.4}{217}{217}
\problemtext 无限长半径$R$的实心圆柱导体均匀载流，
在圆柱内与圆柱外，$B$分别如何依赖轴距$r$？
\workings
由横截面电流比例$I_{\mathrm{enc}}=I r^2/R^2$，
内区$B(2\pi r)=\mu_0I_{\mathrm{enc}}$，得$B\propto r$。
外区$I_{\mathrm{enc}}=I$，得$B\propto1/r$。
\[
\boxed{B_{\mathrm{inside}}\propto r,\quad B_{\mathrm{outside}}\propto1/r}.
\]
原书选\textbf{B}。
\crosscheck 不能把导体内场强误写为$1/r$，否则在$r=0$发散。

\sourceproblem{8.1.5}{217}{217}
\problemtext 一根电阻均匀、半径$R$的圆环，在同一直径的$a,c$两点与两根相切的长直导线连接，
两长直导线均向下通电，电流经圆环分流。判断环心$O$的合磁场是否为零。
\workings
两个长直线在圆心两侧到$O$距离相同、方向相反，故
$\bm B_1+\bm B_2=0$。
圆环中从$a$到$c$两段半圆电阻相等、分流各$I/2$；
它们绕$O$方向相反，圆心磁场亦相消，$\bm B_3=0$。
因此
\[
\boxed{\bm B_O=\bm0}.
\]
原书选\textbf{B}。
\crosscheck 两直线对$O$的磁场分别不为零，
而是矢量叠加相消；不能说“直导线不产生磁场”。

\sourceproblem{8.1.6}{217}{217}
\problemtext 一圆形载流导线的平面内，选一个位于其内侧的同心闭合圆周$L$，
应用安培环路定理可推出什么结论？
\workings
$L$围成的平面区域没有被导线穿越，所以
\[
\boxed{\oint_L\bm B\cdot d\bm l=0}
\]
（真空）。
但通电圆环在圆内各点仍可产生非零且一般不沿$L$切向的磁场，
所以不能推出$B=0$。原书选\textbf{B}。
\crosscheck $L$位于载流圆环所在的平面，圆环的磁场常沿平面法向，
与平面内的线元正交，环流自然为零。

\sourceproblem{8.1.7}{217}{217}
\problemtext 通电直线框一段长度$\ell$进入均匀水平磁场，安培力使天平失衡，
须在左盘加质量$m$才能补偿。若磁场增强至三倍、电流减半，
其他条件与方向不变，需要增加多少质量？
\workings
安培力幅值$F=I\ell B$（电流线段与磁场垂直），
故由$mg=I\ell B$得到新需加质量
\[
m'=\frac{(I/2)\ell(3B)}{g}=\boxed{\frac32m}.
\]
原书选\textbf{B}。
\crosscheck 尽管磁场增强，电流同时减少，净比例为$3/2$，
不能只按磁场的三倍判断。

\sourceproblem{8.1.8}{217}{217}
\problemtext 均匀磁场在纸面内向左，矩形电流线圈所在平面与磁场平行；
图中上边电流$d\to a$向左、右边$c\to d$向上，
判定受力后初始转动时的边界运动方向。
\workings
用右手坐标系$x$向右、$y$向上、$z$由纸面向外。
设$\bm B=-B\bm e_x$。
左边$a\to b$电流向下，受力方向
$(-\bm e_y)\times(-\bm e_x)=-\bm e_z$，
即\emph{转入纸内}；右边$c\to d$电流向上，
$(+\bm e_y)\times(-\bm e_x)=+\bm e_z$，即\emph{转出纸面}。
故\textbf{$ab$边转入纸面、$cd$边转出纸面}，对应选项A。
原书选\textbf{A}。
\crosscheck 线圈总平移磁力可为零，但左右两边受到等大反向力，
形成非零力偶而使其转动。

\sourceproblem{8.1.9}{217--218}{217--218}
\problemtext 同一均匀磁场中两个平面通电线圈面积满足$S_1=2S_2$，
电流满足$I_1=2I_2$。求最大磁力矩比$M_1/M_2$。
\workings
线圈转矩大小$M=ISB\sin\theta$（题设单匝或同匝数），
最大值$M_{\max}=ISB$。因此
\[
\boxed{\frac{M_1}{M_2}=\frac{I_1S_1}{I_2S_2}=2\cdot2=4}.
\]
原书选\textbf{C}。
\crosscheck 若两线圈的匝数不同，还需乘$N_1/N_2$；
题面默认两者对应的匝数条件相同。

\sourceproblem{8.1.10}{218}{218}
\problemtext 100匝圆形线圈，半径$5\,\mathrm{cm}$，
电流$0.1\,\mathrm A$，均匀磁场$1.5\,\mathrm T$。
线圈磁矩与磁场方向夹角从$0^\circ$变为$180^\circ$，
求磁场力对线圈所做的功，讨论原书选项给出的正数。
\workings
磁矩大小$\mu=NIS=100(0.1)\pi(0.05)^2$。
磁偶极子势能$U=-\mu B\cos\theta$，
所以\emph{磁场力做功}等于势能的负改变量：
\[
\boxed{W_{\mathrm B}
=-\Delta U=\mu B(\cos180^\circ-\cos0^\circ)
=-2NISB\approx-0.236\,\mathrm J}.
\]
外力若缓慢地将线圈从平衡取向翻转，则需做正功
$\boxed{W_{\mathrm{ext}}=+0.236\,\mathrm J}$。
原书给出的选项均为正数，其中\textbf{A：0.24 J}对应\emph{功的绝对值或外力做功}，而不对应题目逐字所称“磁场力做功”的有符号数值。
\crosscheck \textbf{原题措辞/答案符号待独立二校}。
从稳定平衡到反平衡，磁场做功应为负，不能为了匹配正选项隐去负号。

\sourceproblem{8.1.11}{218}{218}
\problemtext 带电粒子以速率$v$、与磁场夹角$\theta$射入均匀磁场，轨迹为螺旋线。
分别考察改变$B,\theta,v$，哪种变化能增大螺距？
\workings
垂直分速度$v_\perp=v\sin\theta$给圆周运动，
平行分速度$v_\parallel=v\cos\theta$保持恒定。
回旋周期$T=2\pi m/(|q|B)$，
因此
\[
\boxed{h=v_\parallel T
=\frac{2\pi mv\cos\theta}{|q|B}}.
\]
保持其他量不变，减小$B$使螺距增大，原题选\textbf{B}。
\crosscheck 在$0<\theta<90^\circ$中增大$\theta$会减小$\cos\theta$；
增大$B$或减小$v$也会减小螺距。

\sourceproblem{8.1.12}{218}{218}
\problemtext 长直螺线管紧密绕线，长$\ell\gg a$，总匝数$N$，
通电$I$，内填满相对磁导率$\mu_r$的均匀线性磁介质。
问管内磁场强度$H$与磁感应强度$B$。
\workings
取长螺线管内部与外部的矩形安培环路，
$\oint\bm H\cdot d\bm l=N_{\mathrm{enc}}I$，忽略管外场，
\[
\boxed{H=\frac{NI}{\ell}},\qquad
\boxed{B=\mu_0\mu_r\frac{NI}{\ell}}.
\]
在四个原书选项中只有\textbf{D}（$H=NI/\ell$）正确。
\crosscheck $H$主要由自由通电匝数决定；
$\mu_r$应进入$B=\mu H$而不应凭空进入安培环路定理右端。

\sourceproblem{8.1.13}{218}{218}
\problemtext 按相对磁导率$\mu_r$区分顺磁、抗磁、铁磁物质，判断其大小关系。
\workings
通常弱磁场下各向同性材料，
顺磁质$\mu_r$略大于1，
抗磁质$\mu_r$略小于1，
铁磁质通常$\mu_r\gg1$且随磁化状态和外场改变。
故对应
\[
\boxed{\mu_{r,\mathrm{para}}>1,\quad
\mu_{r,\mathrm{dia}}<1,\quad
\mu_{r,\mathrm{ferro}}\gg1}.
\]
原书选\textbf{C}。
\crosscheck 不应把铁磁体的$\mu_r$当作任意条件下不变的常数；
强磁场饱和及磁滞时线性模型不成立。
% === END textbook-ch08-choices.tex ===

% === BEGIN textbook-ch08-blanks.tex blob 2cd55f2dfb013964b5a2abbbd9dcd08a591b65f0 ===
% 《大学物理简明教程》第5版，第8章稳恒磁场填空题8.2.1—8.2.12
% 原书印刷与扫描PDF第218—219页。逐题简述与独立计算；图形歧义标注二校。
\subsection{填空题：磁场环流、安培力与螺旋运动}

\sourceproblem{8.2.1}{218}{218}
\problemtext 有限长载流直导线的磁场可依据哪一基本规律计算？
不具有充分对称性的此类磁场能否仅凭安培环路定理直接计算？
\workings
微元电流产生的磁感应强度满足毕奥--萨伐尔定律，
\[
\boxed{d\bm B=\frac{\mu_0}{4\pi}
\frac{I\,d\bm l\times\hat{\bm r}}{r^2}}.
\]
由于有限直线段没有圆柱平移对称性，
无法简单提取$B$作为环路上的常量，
故仅用安培定理通常\textbf{不能}直接求出局部场。
\crosscheck 安培环路定理本身仍成立，
问题在于缺乏从环流反推出各点场强的足够对称性。

\sourceproblem{8.2.2}{218}{218}
\problemtext 两个尺寸和绕法相同的螺线管分别有或无铁芯，
若通以相等自由电流，比较内部$\bm H$与$\bm B$。
\workings
忽略退磁因子和漏磁，长螺线管
\[
\boxed{H=nI}\quad(n=N/\ell)
\]
仅由自由电流决定，所以两个螺线管内$H$分布\textbf{相同}。
但$B=\mu_0\mu_rH$，
有铁芯者$\mu_r$通常远大于1，所以$B$分布\textbf{不相同}。
\crosscheck 上述近似要求长螺线管、磁芯尚未饱和；
真实有限铁芯的边缘退磁场可使$H$分布有所变化。

\sourceproblem{8.2.3}{218}{218}
\problemtext 原书图8.2.3以三维闭合曲线$L$环绕
多根导线，电流标记为$I_1,I_2$。
写出沿指定取向的$\oint_L\bm B\cdot d\bm l$。
\workings
在真空中安培环路定理精确地给出
\[
\boxed{\oint_L\bm B\cdot d\bm l
=\mu_0I_{\mathrm{link}}^{\mathrm{alg}}},\qquad
I_{\mathrm{link}}^{\mathrm{alg}}=\sum_j\epsilon_j I_j.
\]
$\epsilon_j\in\{-1,0,1\}$为每根导线和定向曲面
相交的带符号链接数，正向由右手法则确定。
对于原图曲面须明确$I_1$的各次穿越是否相互抵消、
$I_2$是否穿越，才能将此式化为具体$I_1,I_2$组合。
\crosscheck \textbf{待高清空间图独立二校：}本稿不猜测三维示意图中的遮挡穿越方向，
仅录入可严谨确证的安培定理与求值规则；
因此本条为\emph{未完成数值化的条件性题解}，不得宣称通过完整人工质检。

\sourceproblem{8.2.4}{218}{218}
\problemtext 半径$R=0.2\,\mathrm m$的上半圆形闭合线圈，
电流$I=5\,\mathrm A$，处于和线圈平面平行的匀强磁场$B=0.5\,\mathrm T$中。
磁场沿直径由$a$指向$c$。
求上半圆弧$abc$所受安培力大小，以及闭合线圈的最大磁力矩。
\workings
匀强场中任意电流导线段受合力
$\bm F=I(\bm r_{\mathrm{end}}-\bm r_{\mathrm{start}})\times\bm B$。
半圆两端间的位移为直径方向，正与$\bm B$平行，故
\[
\boxed{F_{abc}=0}.
\]
半圆面积$S=\pi R^2/2$，磁矩垂直纸面、与磁场夹角$90^\circ$，
\[
\boxed{M=ISB
=5\cdot\frac{\pi(0.2)^2}{2}\cdot0.5
\approx0.157\,\mathrm{N\,m}}.
\]
按图中直径边电流从$a$到$c$，上半圆弧电流从$c$经$b$到$a$，
整圈电流从纸外观察为逆时针，故磁矩指向纸外；
磁场向右，则由$\bm M=\bm\mu\times\bm B$，
\textbf{力矩方向沿图中竖直轴$O'O$向上（从$O'$指向$O$）}。
\crosscheck 虽然圆弧受\emph{合力}为零，但闭合线圈仍受非零磁力矩。
这里力矩方向是\emph{轴向矢量}方向，而非线圈整体平移方向。

\sourceproblem{8.2.5}{218}{218}
\problemtext 电荷$q=8.0\times10^{-9}\,\mathrm C$以
$v=3.0\times10^{-5}\,\mathrm{m/s}$做半径
$R=6.0\times10^{-8}\,\mathrm m$的匀速圆周运动。
求轨道中心的磁感应强度及相应轨道磁矩大小。
\workings
等效电流为$I=q/T=qv/(2\pi R)$，圆心场
\[
\boxed{B=\frac{\mu_0I}{2R}
=\frac{\mu_0qv}{4\pi R^2}
\approx6.67\times10^{-6}\,\mathrm T}.
\]
轨道磁矩为
\[
\boxed{\mu=I\pi R^2=\frac12qvR
=7.20\times10^{-21}\,\mathrm{A\,m^2}}.
\]
\crosscheck $B$随$R^{-2}$增大，$\mu$随$R$增大；
这里$q$取正，若电荷为负磁矩方向与运动角动量相反。

\sourceproblem{8.2.6}{218--219}{218--219}
\problemtext 无限长载流直导线在点$P$处弯成半径$R$的整圆，
圆心$O$距原直线为$R$，电流$I$不变。
写出$O$点磁场（圆形回路与剩余长直导线叠加）。
\workings
圆形部分在圆心的磁场大小为
\[
B_{\mathrm{circle}}=\frac{\mu_0I}{2R}.
\]
长直切线部分在距离$R$处贡献
\[
B_{\mathrm{straight}}=\frac{\mu_0I}{2\pi R}.
\]
二者都沿纸面法线，方向由图中的圆电流绕向和直线电流
右手法则判定，故一般有
\[
\boxed{B_O=\frac{\mu_0I}{2R}
\left|1\pm\frac1\pi\right|}.
\]
当两者同向取$+$，反向取$-$。
\crosscheck \textbf{图示支路方向待二校：}
圆环是完整一圈还是仅插入导线主干的旁路，必须明确分流与
绕向；本条保留条件性符号，不能将其视为确定唯一数值答案。

\sourceproblem{8.2.7}{219}{219}
\problemtext 磁场中一圆形试验线圈面积增为两倍，
通电电流也增为两倍，问在同一磁场中最大磁力矩增为原来的几倍。
\workings
固定匝数$N$的平面线圈磁矩$\mu=NIS$，
最大磁力矩$M_{\max}=NISB$，所以
\[
\boxed{\frac{M'_{\max}}{M_{\max}}
=\frac{(2I)(2S)}{IS}=4}.
\]
\crosscheck 若题面说的是\emph{面积}加倍，不是半径加倍，
故面积比例恰为2；半径加倍才使面积加倍4倍。

\sourceproblem{8.2.8}{219}{219}
\problemtext 氢原子中电子沿半径$r$的圆周运动，
外加匀强磁场$B$与轨道平面平行。
写出等效环电流的最大磁力矩，以电子质量$m_e$、电荷绝对值$e$表示。
\workings
设轨道速率$v$，等效电流大小$I=ev/(2\pi r)$，
轨道面积$\pi r^2$，所以
\[
\boxed{M=\mu B=I\pi r^2B
=\frac{evr}{2}B
=\frac{eL}{2m_e}B},
\qquad L=m_evr.
\]
若另加\emph{玻尔最小角动量轨道}条件$L=\hbar$，
才可简化为$\boxed{M=e\hbar B/(2m_e)}$。
\crosscheck 原文未给电子速率和角动量量子数，
因此仅由$m_e,e,r,B$并不能决定唯一数值；
假定玻尔基态须明确为附加模型，不可从题面凭空推出。

\sourceproblem{8.2.9}{219}{219}
\problemtext 边长$a$、匝数$N$的正方形电流线圈可绕过
相对两边中点的轴自由转动，通电$I$，
与匀强磁场$B$的\emph{磁矩}夹角为$\theta$，
转动惯量$J$，求所受磁转矩。
\workings
线圈磁矩大小$\mu=NIa^2$，电磁转矩
\[
\boxed{M=NIa^2B\sin\theta},\qquad
\bm M=\bm\mu\times\bm B.
\]
动力学角加速度可再由$J\ddot\phi=M$求得，
但题目只问\emph{转矩}，无需使用$J$。
\crosscheck 如果所给$\theta$为\emph{线圈平面}与磁场夹角，
应改写为$M=NIa^2B|\cos\theta|$；
本稿依原文“磁矩与磁场夹角”使用正弦。

\sourceproblem{8.2.10}{219}{219}
\problemtext 长直螺线管单位长度$n$匝，每匝通电$I$，
管内充满相对磁导率$\mu_r$的均匀介质。
求靠近中部的$B,H$。
\workings
安培环路定理作用于磁场强度$\bm H$，
$H\ell=nI\ell$，故
\[
\boxed{H=nI\quad(\mathrm{A/m})},\qquad
\boxed{B=\mu_0\mu_r nI\quad(\mathrm T)}.
\]
\crosscheck 这里$n$的单位是匝数每米而非总匝数，
故无需除以线圈长度第二次。

\sourceproblem{8.2.11}{219}{219}
\problemtext 质量$m$、电量$q$的粒子以速率$v$与匀强磁场
$B$成夹角$\theta$进入，轨迹为螺旋线，求回旋半径、周期和螺距。
\workings
把速度分解为$v_\perp=v\sin\theta$与
$v_\parallel=v\cos\theta$。
洛伦兹力提供向心加速度：
\[
\boxed{r=\frac{mv|\sin\theta|}{|q|B}},\qquad
\boxed{T=\frac{2\pi m}{|q|B}},
\]
\[
\boxed{h=v_\parallel T=
\frac{2\pi mv\cos\theta}{|q|B}}.
\]
\crosscheck 带电粒子运动方向与磁场平行时$r=0$，
不发生回旋；周期表达式是其垂直扰动的特征回旋周期。

\sourceproblem{8.2.12}{219}{219}
\problemtext 原图8.2.12给三条材料的$B$--$H$曲线
I（非线性、逐渐饱和）、II（比真空斜率大的线性直线）、
III（比真空斜率小的线性直线），虚线为真空$B=\mu_0H$。
判断三种材料类型。
\workings
线性材料的斜率为$\mu_0\mu_r$，
大于真空$\mu_0$者为顺磁质，小于者为抗磁质。
非线性、磁导率大并逐渐进入饱和的曲线是铁磁质。
因此
故\textbf{曲线II为顺磁质、曲线III为抗磁质、曲线I为铁磁质}。
\crosscheck 铁磁体$B$--$H$曲线的曲率不能用于估计单一常数$\mu_r$，
磁滞回线需另给加载历史。
% === END textbook-ch08-blanks.tex ===

% === BEGIN textbook-ch08-solutions-001-007.tex blob 6c09c4e14810ac4528b340848f9294dccc74d7fe ===
% 《大学物理简明教程》第5版第8章，解答题8.3.1—8.3.7。
% 源扫描PDF及印刷第219—220页；保留原题号/模型，独立LaTeX推导、必要时标注条件。

\subsection{解答题：电流单位、磁通量与有限导体场}

\sourceproblem{8.3.1}{219}{219}
\problemtext 手机电池容量常以毫安时（mAh）表示，
解释该单位的物理意义并将其换算为国际单位制。
\workings
电池容量的$mAh$是\emph{电荷量}而不是能量或功率。
若输出电流恒定，所通过的总电量
$Q=It$，因$1\,\mathrm{h}=3600\,\mathrm s$，
\[
\boxed{1\,\mathrm{mAh}=10^{-3}\,\mathrm A\cdot3600\,\mathrm s
=3.6\,\mathrm C}.
\]
例如标称$3000\,\mathrm{mAh}$对应约$1.08\times10^4$ C的标称电荷量。
\crosscheck 若需要电池能量还须给电压：
$E\simeq UQ$，不能把mAh本身直接说成焦耳。

\sourceproblem{8.3.2}{219}{219}
\problemtext 磁感线上各点$B$的\emph{数值}是否相同？
为何不能把对运动电荷的磁场力方向定义为磁感应强度的方向？
\workings
磁感线的切线表示局部$\bm B$的\emph{方向}，
但磁感线不同位置疏密可变，所以$B$的\emph{大小}一般不恒定。
对带电粒子，洛伦兹磁力为
\[
\boxed{\bm F_{\mathrm B}=q\bm v\times\bm B}.
\]
因此$\bm F_{\mathrm B}$同时垂直于$\bm v,\bm B$，
在任意非零速度情况下都不平行于$\bm B$。
更重要的是$\bm F$还取决于$\bm v$与$q$的符号，
不能用某一次受力方向定义$\bm B$。
\crosscheck $\bm v\parallel\bm B$时磁力为零，
但所在点可以存在非零磁场。

\sourceproblem{8.3.3}{219}{219}
\problemtext 对一段\emph{有限长}稳恒载流直导线，
能否只凭安培环路定理求出其周围各点的磁场？
\workings
一般的安培环路定理为
\[
\oint_L\bm B\cdot d\bm l=\mu_0I_{\mathrm{enc}}^{\mathrm{alg}}
\quad(\text{SI units}).
\]
有限导线缺乏无限长直导线的平移及圆柱对称性，
即使知道环路上的积分值，也不能因此令$B$在整条环路上大小相同。
故\textbf{不能仅依此定理直接求场强分布}。
相应有限直线段的场可利用毕奥--萨伐尔定律积分：
\[
\boxed{B=\frac{\mu_0I}{4\pi r}
(\sin\alpha_1+\sin\alpha_2)}
\]
（$\alpha_i$为两端点从垂足方向量得的对应角，
且观察点到导线延长线的垂距为$r$）。
\crosscheck 有限孤立直线段不能单独构成完整稳恒闭合电路，
实际回路中的返流也可能影响外部磁场；上述公式仅是该导线段的贡献。

\sourceproblem{8.3.4}{219}{219}
\problemtext 原图8.3.4给出三维框架$abcdef$，均匀磁场
$\bm B=2.0\,\mathrm T\,\bm e_x$，有两段长30 cm及沿$x$方向长50 cm。
比较穿过$abcd$、$bcfe$、$aefd$三个面的磁通量。
\workings
磁通量只取决于面积矢量沿磁场方向的投影：
\[
\boxed{\Phi_S=\int_S\bm B\cdot d\bm S
=B\,A_{S,\perp x}}.
\]
按照原图三个顶点坐标的几何关系，
$abcd$处于$yz$平面，垂直磁场，其两边投影长均为$0.30$ m，
故磁通量大小
\[
\boxed{|\Phi_{abcd}|=2.0(0.30)^2=0.180\,\mathrm{Wb}}.
\]
$bcfe$平行$x$轴，面法向与$\bm B$垂直，
所以$\boxed{\Phi_{bcfe}=0}$。
斜面$aefd$在$yz$平面的正交投影与$abcd$相同，
故
\[
\boxed{|\Phi_{aefd}|=0.180\,\mathrm{Wb}}.
\]
\crosscheck \textbf{立体图的法向取向需二校}：
两个非零通量的\emph{代数正负}取决于曲面定向，
本稿只给无歧义的大小；闭合曲面的有向通量总和必须为零。

\sourceproblem{8.3.5}{219}{219}
\problemtext 长直导线经过半径$R$的四分之一圆弧后，
连接两端的半无限长直线，通电$I$，图中圆心为$O$。
求该点磁感应强度。
\workings
按图，左边水平直线与$O$共线，其电流元$d\bm l$
与指向$O$的位矢平行，故贡献为零。
四分之一圆弧张角$\pi/2$，
\[
B_{\mathrm{arc}}=\frac{\mu_0I}{4\pi R}\frac{\pi}{2}
=\frac{\mu_0I}{8R}.
\]
右侧半无限长水平直线距$O$为$R$，
由毕奥--萨伐尔定律
\[
B_{\mathrm{half}}=\frac{\mu_0I}{4\pi R}.
\]
依原图电流方向，圆弧和右侧直线在$O$处都指向纸内，故
\[
\boxed{B_O=\frac{\mu_0I}{8R}
+\frac{\mu_0I}{4\pi R}
=\frac{\mu_0I}{8\pi R}(\pi+2)}
\]
方向垂直纸面向里。
\crosscheck 若把右边的\emph{半无限长}导线按无限长计算，
会把它的贡献误算为两倍。

\sourceproblem{8.3.6}{219--220}{219--220}
\problemtext 相互平行的无限长导线$L_1,L_2$相距$0.10$ m，
载流分别为$I_1=20$ A向右、$I_2=10$ A向左。
点$A$位于两线之间距$L_2$为$0.05$ m，
点$B$位于$L_2$下方$0.05$ m。
求$A,B$处磁场、以及外部磁场为零的位置。
\workings
单根无限长直线在距线$r$处
$B=\mu_0I/(2\pi r)$。
在$A$处，距两导线分别为$0.05$ m，磁场均向纸内，
\[
\boxed{B_A=\frac{\mu_0}{2\pi}
\left(\frac{20}{0.05}+\frac{10}{0.05}\right)
=1.20\times10^{-4}\,\mathrm T}
\]
方向向纸内。
在$B$处距上导线为$0.15$ m、距下导线为$0.05$ m，
下线磁场向外，上线磁场向里，故
\[
\boxed{B_B=\frac{\mu_0}{2\pi}
\left(\frac{10}{0.05}-\frac{20}{0.15}\right)
\approx1.33\times10^{-5}\,\mathrm T}
\]
方向向纸外。
令$L_2$下方距其$r$的磁场合成零，
\[
\frac{10}{r}=\frac{20}{r+0.10}
\quad\Longrightarrow\quad
\boxed{r=0.10\,\mathrm m}.
\]
\crosscheck 两线之间电流相反导致磁场同向相加，
零场点只能在电流较小的$L_2$一侧外部。

\sourceproblem{8.3.7}{219--220}{219--220}
\problemtext 粗细均匀、各段电阻与弧长成正比的导电圆环，
在$A,B$两点接上两条沿过圆心直线方向延伸的导线，电流从$A$流入、从$B$流出。
求圆心$O$处磁感应强度。
\workings
从$A$到$B$有两条圆弧支路，其圆心角分别$\theta$与$2\pi-\theta$。
电阻比与弧长比相同：
$R_1/R_2=\theta/(2\pi-\theta)$，
分流电流满足$I_1/I_2=R_2/R_1$，
故$I_1\theta=I_2(2\pi-\theta)$。
弧长产生圆心场
$B_1=\mu_0I_1\theta/(4\pi R)$、
$B_2=\mu_0I_2(2\pi-\theta)/(4\pi R)$，
两支路绕向相反，恰好抵消。
两根直线连接线沿$O$的径向，
其每个电流元和指向$O$的矢量平行，不产生圆心磁场。
因此
\[
\boxed{\bm B_O=\bm0}.
\]
\crosscheck 平衡所需的不是两支路电流相等，
而是$I_1\theta=I_2(2\pi-\theta)$，
对应中心场等大反向。
% === END textbook-ch08-solutions-001-007.tex ===

% === BEGIN textbook-ch08-solutions-008-014.tex blob e91a5e269b143915f033e47b97aecccfc16c820d ===
% 《大学物理简明教程》第5版第8章，解答题8.3.8—8.3.14。
% 源扫描及印刷第220—221页。三维图和电磁弹射模型注明条件，待独立二校。

\subsection{解答题：环路定理、安培力和霍尔电压}

\sourceproblem{8.3.8}{220}{220}
\problemtext 原图8.3.8中两条穿出纸面的长直导线，
左线电流$I_2$指向纸内，右线电流$I_1$指向纸外。
考察分别围住左线、右线及两线的三条闭合曲线$a,b,c$：
求各环路的$\oint\bm B\cdot d\bm l$，
并判断各曲线上磁场是否处处相同或处处为零。
\workings
令逆时针为每条闭合曲线的正向（相应面法向向纸外）。
安培环路定理给
\[
\boxed{\oint_a\bm B\cdot d\bm l=-\mu_0I_2},\qquad
\boxed{\oint_b\bm B\cdot d\bm l=+\mu_0I_1},
\]
\[
\boxed{\oint_c\bm B\cdot d\bm l=\mu_0(I_1-I_2)}.
\]
如果图题指定的环路正向与以上相反，相应积分整体变号。
即使$|I_1|=|I_2|$使最外圈总环流为零，
磁场一般仍由两导线叠加，\emph{不能}推出各点$B=0$；
且非同轴环路上大小通常也不是常数。
\crosscheck 三种闭合路径的源电流是\emph{代数和}，
不可将它们全部取绝对值再相加。

\sourceproblem{8.3.9}{220}{220}
\problemtext 无限长、内半径$a$、外半径$b$的空心圆柱导体，
电流$I$沿轴向且均匀分布于导体截面环形区域。
证明在$a<r<b$处的磁感应强度公式。
\workings
体电流密度
$J=I/[\pi(b^2-a^2)]$。
取半径$r$的同轴圆形安培环路，
包围自由电流为
\[
I_{\mathrm{enc}}=J\pi(r^2-a^2)
=I\frac{r^2-a^2}{b^2-a^2}.
\]
由$B(r)(2\pi r)=\mu_0I_{\mathrm{enc}}$，
\[
\boxed{B(r)=\frac{\mu_0I}{2\pi(b^2-a^2)}
\frac{r^2-a^2}{r}},\quad a<r<b.
\]
完整分段还有$B=0$（$r<a$）和
$B=\mu_0I/(2\pi r)$（$r>b$）。
\crosscheck 在$r=a$处场强为零，在$r=b$处连续接到无限长直导线外场公式。

\sourceproblem{8.3.10}{220}{220}
\problemtext 真空中一无限长实心圆柱导体半径$R$，
电流$I$均匀分布于截面内。
图8.3.10阴影为与圆柱轴平行、轴向长$L=1.0\,\mathrm m$、
在径向从$r=0$延伸到$r=2R$的面。
求穿过该阴影面的磁通量。
\workings
均匀体电流的磁场方位角方向，
在阴影径向—轴向平面上恰垂直该面。
由安培定理
\[
B(r)=
\begin{cases}
\mu_0Ir/(2\pi R^2),&r<R,\\
\mu_0I/(2\pi r),&r\ge R.
\end{cases}
\]
故取沿该方位磁场的有向面法线，
\[
\Phi=L\int_0^{2R}B(r)\,dr
=\frac{\mu_0IL}{2\pi}
\left[\int_0^R\frac r{R^2}\,dr+
\int_R^{2R}\frac{dr}{r}\right].
\]
得到
\[
\boxed{\Phi=\frac{\mu_0IL}{2\pi}
\left(\frac12+\ln2\right)}.
\]
\crosscheck 此处图示阴影从\emph{轴线}延至$2R$，
故必须包括圆柱内部的$1/2$项；
若仅取圆柱外部$R\to2R$则只有$\ln2$项。

\sourceproblem{8.3.11}{220}{220}
\problemtext 长同轴电缆由实心内导体半径$a$和空心外圆筒导体内、外半径$b,c$组成，
电流$I$从内导体流出并由外筒以$-I$返回，两者各自均匀分布在横截面中。
求各半径区域的磁场强度大小。
\workings
对半径$r$的同轴圆形安培环路，
根据环路包围的代数电流不同逐段计算：
\[
I_{\mathrm{enc}}(r)=
\begin{cases}
I r^2/a^2,&0\le r<a,\\
I,&a\le r<b,\\
I\left[1-\frac{r^2-b^2}{c^2-b^2}\right],&b\le r<c,\\
0,&r\ge c.
\end{cases}
\]
故
\[
\boxed{B(r)=\frac{\mu_0}{2\pi r}
\begin{cases}
I r^2/a^2,&r<a,\\
I,&a<r<b,\\
I(c^2-r^2)/(c^2-b^2),&b<r<c,\\
0,&r>c.
\end{cases}}
\]
方位由内导体电流的右手法则确定。
\crosscheck $r=a,b,c$的相邻表达式连续；
$r>c$时包围$+I$和$-I$相消，外部场为零。

\sourceproblem{8.3.12}{220--221}{220--221}
\problemtext 图8.3.12以磁感应强度$B=0.05\,\mathrm T$的均匀横向场
测量进入纸面的血液流速。
血管直径$d=3\,\mathrm{mm}$，
上下电极间霍尔电压绝对值$U_H=150\,\mu\mathrm V$；
图示上极板为负、下极板为正。
求血流速度，并判断净漂移载流子的电荷符号。
\workings
在稳定的横向电荷分离中，电力与磁力相互平衡：
$|q|E=|q|vB$，$E=U_H/d$，因此
\[
\boxed{v=\frac{U_H}{Bd}
=\frac{150\times10^{-6}}{(0.05)(3\times10^{-3})}
=1.00\,\mathrm{m/s}}.
\]
图示磁场从$N$极指向$S$极，朝右；流速指向纸内。
取$x$向右、$y$向上、$z$向纸外，
$\bm v=-v\bm e_z,\bm B=B\bm e_x$，
则$\bm v\times\bm B=-vB\bm e_y$，即对正电荷作用力朝下。
下电极为正、上电极为负，故按图示的简单单载流子模型，
迁移电荷为\textbf{正电荷}。
\crosscheck 血液是多离子电解质，真实电极的电化学条件复杂；
本题按均匀体电磁流量计的简化霍尔模型处理。

\sourceproblem{8.3.13}{221}{221}
\problemtext 图8.3.13两条平行金属导轨长$L$、间距$d$、
每条导轨圆柱半径$R$，$L\gg d$。
质量$m$的导电弹射体搭接两轨，以恒定电流$I$加速，
从轨中点出发、再前进$L/2$离开。估算安培力及出轨速率。
\workings
电流经一根导轨、弹射体横杆再经另一导轨反向返回，
在横杆处形成横向磁场。
\emph{若将两导轨在弹射体附近近似为无限长直导线}，
横杆上距第一轨轴线$y$的磁感应强度
\[
B(y)\approx\frac{\mu_0I}{2\pi}
\left(\frac1y+\frac1{d-y}\right).
\]
用导轨半径作为积分截止尺度，从$y=R$至$d-R$积分，得
\[
\boxed{F_{\mathrm{long}}=
I\int_R^{d-R}B(y)\,dy
=\frac{\mu_0I^2}{\pi}
\ln\frac{d-R}{R}}.
\]
假定此力沿出轨方向恒定、忽略摩擦与空气阻力，
由功—能定理$F(L/2)=\tfrac12mv^2$，
\[
\boxed{v_{\mathrm{exit}}\approx
\sqrt{\frac{F_{\mathrm{long}}L}{m}}}.
\]
\crosscheck \textbf{边界模型待二校}：若两导轨只向弹射体一侧延伸、
各在横杆位置表现为\emph{半无限长}电流，则每轨的场需近似减半，
$F$和$v$相应改变；由于题给$R$且无实际电源与端部回路详细几何，
不得将此无限直线近似宣称成唯一精确结果。

\sourceproblem{8.3.14}{221}{221}
\problemtext 图8.3.14一根电流为$I$的弯曲导线位于均匀磁场中，
曲线两端为$a,b$，磁场$B$垂直纸面向里。
求整段导线所受安培合力，讨论导线形状是否影响结果。
\workings
安培力微元$d\bm F=I\,d\bm l\times\bm B$，
均匀磁场可移出线积分，
\[
\bm F=I\int_a^b d\bm l\times\bm B
=I\left(\int_a^b d\bm l\right)\times\bm B.
\]
位移积分$\int_a^b d\bm l=\bm r_b-\bm r_a$，
故
\[
\boxed{\bm F=I(\bm r_b-\bm r_a)\times\bm B}.
\]
若端点连线在纸面内长为$\ell_{ab}$，
由于$\bm B$垂直纸面，合力大小为
$\boxed{F=IB\ell_{ab}}$，方向由右手法则判定。
\crosscheck $\ell_{ab}$是两端点的\emph{直线距离}，
不是弯曲导线的弧长；
图中受力矢量方向取决于电流从$a$到$b$还是相反。
% === END textbook-ch08-solutions-008-014.tex ===

% === BEGIN textbook-ch08-solutions-015-021.tex blob f77788111e4dbbb47d2b7b22a3af58f2a3d22ff4 ===
% 《大学物理简明教程》第5版第8章解答题8.3.15—8.3.21，
% 原书与扫描PDF第221—222页。源条件逐页校对，原创推导，保留图示疑点。
\subsection{解答题：磁力、回旋粒子与磁性材料}

\sourceproblem{8.3.15}{221}{221}
\problemtext 长直导线$AB$沿竖直方向通过电流$I_1=20$ A，
右侧与其共面的矩形线圈$CDEF$通电$I_2=10$ A。
线圈高$b=20.0$ cm、宽$a=9.0$ cm，
最近边$CD$与长直线距离$d=1.0$ cm。
按原图线圈电流顺时针，求四边受力、总合力和合力矩。
\workings
导线右侧磁场大小$B(r)=\mu_0I_1/(2\pi r)$，
方向指向纸内。
与长直线平行的近边$CD$电流向上，
磁力朝长直线，大小
\[
\boxed{F_{CD}=\frac{\mu_0I_1I_2b}{2\pi d}
=8.00\times10^{-4}\,\mathrm N}.
\]
远边$EF$电流向下，受力背离长直线，大小
\[
\boxed{F_{EF}=\frac{\mu_0I_1I_2b}{2\pi(d+a)}
=8.00\times10^{-5}\,\mathrm N}.
\]
上边$CF$和下边$DE$分别受向上或向下的力，
其大小同为
\[
\boxed{|F_{CF}|=|F_{DE}|=
\frac{\mu_0I_1I_2}{2\pi}
\ln\frac{d+a}{d}
\approx9.21\times10^{-5}\,\mathrm N},
\]
方向相反而相互抵消。
因此线圈合力
\[
\boxed{F_{\mathrm{net}}=F_{CD}-F_{EF}
=7.20\times10^{-4}\,\mathrm N}
\]
指向长直导线。对线圈中心求各力矩，平行边受力作用线经过
对应边的中心高度，水平边力矩彼此抵消，故
$\boxed{M_{\mathrm{net}}=0}$。
\crosscheck 平面内该回路磁矩与长直导线所产生的$\bm B$
平行或反平行，不存在使其绕平面内轴转动的净力偶；
但仍有朝向长直线的非零合力。

\sourceproblem{8.3.16}{221}{221}
\problemtext 一片厚度$d=1.0\,\mathrm{mm}$的铜片，
磁场$B=1.5\,\mathrm T$垂直于片面。
铜片每立方厘米有$8.4\times10^{22}$个自由电子，
通过垂直于磁场方向的电流$I=200$ A，
求铜片两侧霍尔电势差的大小。
\workings
电子数密度换算
$n=8.4\times10^{22}\,\mathrm{cm^{-3}}
=8.4\times10^{28}\,\mathrm{m^{-3}}$。
在简单单载流子霍尔模型中
\[
|U_H|=\frac{BI}{n e d},
\]
代入$e=1.602176634\times10^{-19}$ C得
\[
\boxed{|U_H|=\frac{(1.5)(200)}
 {(8.4\times10^{28})(1.602176634\times10^{-19})(10^{-3})}
\approx2.23\times10^{-5}\,\mathrm V}
\]
即约$22.3\,\mu\mathrm V$。
\crosscheck 铜以电子为主要导电载流子，霍尔系数为负；
题只问霍尔电势差大小，故不凭缺少电极标记的图去猜电压符号。

\sourceproblem{8.3.17}{222}{222}
\problemtext 无限长竖直导线载有向上的电流$I_1=20$ A。
右侧与其共面的水平直导线段$ab$载流$I_2=10$ A，由近端$a$
距竖线$0.01$ m延伸至远端$b$距竖线$0.10$ m，
电流从$a$向$b$。求$ab$所受总安培力以及它对图中$O$点的磁力矩。
\workings
竖直向上的长直电流在其右侧产生指向纸内的磁场
$B(r)=\mu_0I_1/(2\pi r)$。
水平电流段元$d\bm l=dr\,\bm e_x$，
故$d\bm F=I_2d\bm l\times\bm B$
沿纸面\emph{向上}，大小
\[
\boxed{F_{ab}=\frac{\mu_0I_1I_2}{2\pi}
\int_{0.01}^{0.10}\frac{dr}{r}
=\frac{\mu_0I_1I_2}{2\pi}\ln10
\approx9.21\times10^{-5}\,\mathrm N}.
\]
对位于竖直线上的$O$求$z$向力矩，
$dM_O=r\,dF$，所以
\[
\boxed{M_O=\frac{\mu_0I_1I_2}{2\pi}
\int_{0.01}^{0.10}dr
=3.60\times10^{-6}\,\mathrm{N\,m}}.
\]
方向绕$O$逆时针，力矩向纸外。
\crosscheck 力的积分随距离呈对数关系，
力矩积分因力臂乘上$1/r$恰消去距离依赖。

\sourceproblem{8.3.18}{222}{222}
\problemtext 电子在磁感应强度$B=70\times10^{-4}$ T
的均匀磁场中做半径$r=3.0$ cm的圆周运动，
磁场垂直纸面向外。图中电子在$A$点的速度方向向上。
求回旋方向、速率和动能。
\workings
电子电荷为$-e$。在$A$点，
$\bm v=+v\bm e_y$，$\bm B=+B\bm e_z$，
$\bm v\times\bm B$指向$+x$，再乘$-e$则力指向$-x$，
所以圆心位于$A$的左侧、轨迹在纸面内\textbf{逆时针}。
洛伦兹力提供向心力：
\[
evB=\frac{m_ev^2}{r}
\quad\Longrightarrow\quad
\boxed{v=\frac{eBr}{m_e}
\approx3.69\times10^7\,\mathrm{m/s}}.
\]
动能
\[
\boxed{E_k=\frac12m_ev^2
=\frac{e^2B^2r^2}{2m_e}
\approx6.20\times10^{-16}\,\mathrm J
\approx3.87\,\mathrm{keV}}.
\]
\crosscheck $v/c\approx0.12$，教材使用的非相对论公式为近似；
若要求高精度，应按相对论动量$p=\gamma m_ev$修正。

\sourceproblem{8.3.19}{222}{222}
\problemtext 电子在$B=20\times10^{-4}$ T的均匀磁场中沿螺旋线运动，
回旋半径$R=2.0$ cm，螺距$h=5.0$ cm。
(1) 求总速率；(2) 依据图示回旋方向解释磁场轴向的正反判断。
\workings
圆周运动的垂直速率与回旋周期为
\[
v_\perp=\frac{eBR}{m_e}
\approx7.04\times10^6\,\mathrm{m/s},
\qquad
T=\frac{2\pi m_e}{eB}
\approx1.79\times10^{-8}\,\mathrm s.
\]
沿场方向的速率大小满足
$v_\parallel=|h|/T$，所以
\[
v_\parallel\approx2.80\times10^6\,\mathrm{m/s},
\qquad
\boxed{v=\sqrt{v_\perp^2+v_\parallel^2}
\approx7.58\times10^6\,\mathrm{m/s}}.
\]
电子带负电，回旋角速度矢量与$\bm B$的方向反向；
观察原图绕轴旋转的朝向，使用右手定则判出角速度矢量，
磁场再取其反方向。
\crosscheck \textbf{图示轴向二校：}速率和周期不依赖磁场的正负方向；
但确定$\bm B$向上还是向下必须区分从哪一端观察螺旋旋向，
不能仅据“螺旋向上”推断磁场向上。

\sourceproblem{8.3.20}{222}{222}
\problemtext 螺线管平均中心周长$L=0.10$ m，总匝数$N=200$，
通电$I=0.100$ A。
(1) 空气芯求$H,B_0$；
(2) 改为相对磁导率$\mu_r=4200$的均匀线性磁芯求$H,B$。
\workings
使用题设\emph{平均磁路长度}$L$，
\[
\boxed{H=\frac{NI}{L}
=\frac{200(0.100)}{0.10}
=200\,\mathrm{A/m}}.
\]
空气芯近似$\mu_r=1$，
\[
\boxed{B_0=\mu_0H\approx2.51\times10^{-4}\,\mathrm T}.
\]
对线性磁性材料
\[
\boxed{B=\mu_0\mu_rH
=4200B_0\approx1.06\,\mathrm T},\qquad
\boxed{H=200\,\mathrm{A/m}}.
\]
\crosscheck 磁芯改变了$B$，但在理想环形磁路、固定自由电流下
$H$由安培定理保持不变；
铁芯接近饱和时常数$\mu_r=4200$不再适用。

\sourceproblem{8.3.21}{222}{222}
\problemtext 为测环形铁芯的相对磁导率$\mu_r$，
将横截面积$S=0.50\times10^{-4}\,\mathrm{m^2}$的铁芯绕成
平均周长$L=0.10$ m、总匝数$N=200$的环形螺线管。
电流$I=0.10$ A时穿过环芯截面的磁通量为
$\Phi=6.0\times10^{-5}$ Wb，求$\mu_r$。
\workings
磁感应强度近似均匀，
\[
\boxed{B=\frac{\Phi}{S}
=\frac{6.0\times10^{-5}}{0.50\times10^{-4}}
=1.20\,\mathrm T}.
\]
由安培环路定理
\[
\boxed{H=\frac{NI}{L}
=\frac{200(0.10)}{0.10}
=200\,\mathrm{A/m}}.
\]
故有效相对磁导率
\[
\boxed{\mu_r=\frac{B}{\mu_0H}
\approx\frac{1.20}{(4\pi\times10^{-7})(200)}
\approx4.77\times10^3}.
\]
\crosscheck 此为在指定$B,H$工作点测得的\emph{有效割线磁导率}，
铁磁体通常非线性且有磁滞，不能断言在任意$H$下$\mu_r$都相同。
% === END textbook-ch08-solutions-015-021.tex ===
